% Manuscript for epic E5 (#7): the 2016 RBF-FD heat-transfer construction
% ported, and the seed stencils for sub-grid smooth edges of
% docs/stiff-diffusion.md, written up for arXiv. Scaffold from E5.1 (#42);
% each section is drafted by the ticket named in its TRACE block.
%
% SOURCE-OF-TRUTH RULE: docs/stiff-diffusion.md (the stiff note) is canonical
% for the stiff-edge study and docs/port-notes.md for the port (plan D11).
% This manuscript QUOTES them; it never becomes a second source of truth.
% Every number in this file carries a TRACE comment naming a notes section
% or a results file under paper/data/ (E5.3, #44), and the assembly pass
% (E5.10, #51) re-checks all of them.
%
% Section map (plan section 6, E5; the stiff note's section 5.7 for section 6):
%   Sec. 1  Introduction              stiff 1.1, 2.5, 5.3; LITERATURE.md 6a  #45
%   Sec. 2  Setting                   stiff 1.1, 3.1; port notes 1-2          #45
%   Sec. 3  Seeds in one dimension    stiff 1.2-1.8                           #46
%   Sec. 4  1-D results               stiff 2.1-2.5                           #47
%   Sec. 5  Seeds in two dimensions   stiff 3.2-3.4, 3.10, 3.11, 4.10, 5.2    #48
%   Sec. 6  2-D results               stiff 4.1-4.10, 5.1, 5.3                #49
%   Sec. 7  Discussion                stiff 2.5, 5.5-5.6; LITERATURE.md 6b    #50
%   Literature #43 . Figures, tables, numbers #44 . Assembly #51 . arXiv #52
%
% Build: tectonic main.tex   (see paper/README.md)
\documentclass[11pt]{amsart}

\usepackage{amsmath,amssymb,amsthm}
\usepackage{graphicx}
\usepackage{booktabs}
\usepackage[colorlinks=true, linkcolor=blue, citecolor=blue, urlcolor=blue]{hyperref}
\graphicspath{{figures/}}
% Long citation keys and verbatim ledger sentences leave a few lines a
% few points too wide; this lets TeX loosen those lines instead.
\setlength{\emergencystretch}{1.5em}
% arXiv builds with pdfTeX, tectonic with XeTeX. Only the bibliography may
% carry non-ASCII input (accented surnames as UTF-8, paper/README.md), which
% pdfTeX reads with this; make_arxiv.py refuses anything else.
\usepackage{iftex}
\ifPDFTeX\usepackage[utf8]{inputenc}\fi

\newtheorem{theorem}{Theorem}[section]
\newtheorem{lemma}[theorem]{Lemma}
\newtheorem{proposition}[theorem]{Proposition}
\newtheorem{corollary}[theorem]{Corollary}
\theoremstyle{definition}
\newtheorem{definition}[theorem]{Definition}
\theoremstyle{remark}
\newtheorem{remark}[theorem]{Remark}

% Recurring quantities. Only macros the current text uses are defined; the
% section drafts add their own as they need them.
\newcommand{\px}{\partial_x}
\newcommand{\py}{\partial_y}
\newcommand{\pt}{\partial_t}
% A run-in paragraph heading.
\newcommand{\para}[1]{\par\smallskip\noindent\textit{#1}.\enspace}
% Scientific notation for quoted measurements: $\sci{6.2}{-8}$.
\newcommand{\sci}[2]{#1 \times 10^{#2}}
% Draft stubs, visible in the PDF until the owning ticket replaces them.
% make_arxiv.py refuses to package while one remains.
\newcommand{\stub}[1]{\par\noindent\textit{[Stub, #1.]}\par}

% Working title; the final wording is E5.10's (#51).
\title[Seed stencils for diffusion]{Seed stencils for diffusion through
material edges too steep for the grid}
\author{Bradley P. Martin}
\email{1bradley.martin1@gmail.com}
% Fixed revision date, updated by hand with content changes. \today would
% make the committed main.pdf non-reproducible from the committed source.
\date{September 24, 2026}

\begin{document}

% ---------------------------------------------------------------------------
% ABSTRACT. First draft by #45 from the statements the notes make; the final
% wording is #51's. Novelty only in LITERATURE.md section 6's words (#43):
% the fourth and fifth sentences are 6b's abstract-safe sentence verbatim, and
% the last two are 6b's required caveat verbatim, since the comparison with
% the coefficient treatments is a headline here.
% TRACE: the knee, stiff 2.5 (1) and 5.3 (2); the construction's floor of
% order delta and its growth once the edge is resolved, stiff 2.5 (1) and
% 5.3 (3); the seeds, stiff 2.5 (2) and 5.3 (4); the caveat's numbers:
%   801 nodes          stiff 2.5 (3), paper/data/heat1d_stiff.json
%   1.75               stiff 5.3 (10), paper/data/heat2d_stiff_treatments*.json
%   2.3 to 4.7 orders over 1250 to 40,000 nodes
%                      stiff 5.3 (10), paper/data/heat2d_stiff_treatments.json
% (scripts/paper_numbers.py: where "abstract, §1; stiff ...").
% ---------------------------------------------------------------------------
\begin{abstract}
A diffusion coefficient can vary smoothly on the scale of the grid, jump
across an interface, or change smoothly over a width $\delta$ smaller than
the node spacing $h$. Through such a sub-grid edge, standard high-order
stencils that sample the coefficient at the nodes are first order until the
grid resolves the edge, with a knee at $h \approx \delta$, and the
interface-aware stencils built for a jump, applied at the edge's centre, sit
on an error floor of order $\delta$ and leave it once the grid resolves the
edge. We keep the nodes and replace the polynomial basis of the stencils
that see the edge by \emph{seeds}. The seeds are the profiles of solutions
polynomial in time, continued through the edge by ordinary differential
equations along the normal or along the edge's own coordinates; in one
dimension they are the classical formal powers of the Sturm--Liouville
operator, and as $\delta \to 0$ they reduce to the interface-aware stencils
of a jump. On scattered nodes the resulting RBF-FD stencils are fourth order
at every $\delta$ tested, elliptic and parabolic, with no switch between
regimes. We compare them with the coefficient treatments of diffusion
practice. The comparison is against coefficient sampling on the same nodes
and against our implementation of the cell harmonic and arithmetic means, a
widened edge and, in one dimension, the conservative scheme with exact face
conductances, on one contrast and one test problem per dimension. Through an
edge the grid does not resolve none of them is better than second order; the
conservative scheme is exact at the one-dimensional equilibrium and, on the
sinusoidal medium at $\delta = 0$, ahead of the seeds below 801 nodes, where
the seeds' line is pre-asymptotic; on the scattered nodes none improves on
sampling by more than a factor of 1.75 in the RMS norm, and the seeds lie
2.3 to 4.7 orders below the best of them over 1250 to 40,000 nodes.
\end{abstract}

\maketitle

% ===========================================================================
\section{Introduction}
\label{sec:intro}
% TRACE: stiff 1.1 (the setting, the three regimes), 1.2 (the chain), 1.7
% and 2.2 (the construction on a smooth edge), 2.5 (what the 1-D study
% states), 5.3 (what the 2-D study states); plan sections 1 and 3;
% LITERATURE.md 6a for the prior-work paragraph and 6b for the claim
% sentences (#43). Drafted by #45. Numbers are traced paragraph by
% paragraph below; every cache-backed one is a check in
% scripts/paper_numbers.py whose `where` starts with "§1" (or "abstract").

Heat conduction through a composite material is governed by
\begin{equation}
  \label{eq:heat}
  u_t = \nabla \cdot (\alpha \nabla u) \equiv L u
\end{equation}
and, at equilibrium, by $L u = 0$, with a conductivity $\alpha > 0$ that
changes sharply between materials. At the two ends of the range of media
the numerical picture is settled. Where $\alpha$ varies smoothly on
the scale of the grid, a high-order stencil samples it at the nodes and
keeps its order. Where $\alpha$ jumps across an interface, the temperature
and the flux $\alpha\, \partial_n u$ are continuous while the gradient is
not, a standard stencil across the interface is first order, and the
remedies rebuild the stencils or the elements that the interface cuts so
that they satisfy the interface conditions (Section~\ref{sec:priorwork}).
One of them, the interface-aware RBF-FD stencils of
\cite{Martin2016,MartinFornberg2017heat} that this paper extends, keeps
fourth order on scattered nodes, elliptic and parabolic, with a few rows of
nodes placed across each interface. Between the two ends lies a third case:
an edge that is smooth, so that no interface condition applies, but
changes over a width $\delta$ smaller than the node spacing $h$. Every
material edge is smooth on some scale, a model that resolves an edge on a
fine grid presents it as sub-grid on a coarser one, and any regularisation
of a jump produces one.

% TRACE: the 1-D knee, stiff 2.5 (1) (the two-constant medium's ramp
% problem at delta = 0.01 and 0.0025, h = 2 delta and h = delta/2; the drop
% "100 to 200 times", checked as bounds); the 2-D knee, stiff 5.3 (2)
% (0.86-0.94 of the jump's error while h >~ 6 delta, 0.081-0.084 at
% h ~ delta, the bulk stencils' order 5 below the knee). "Falls by four" is
% the first-order line over a factor 4 in h, not a measurement.
A standard stencil through a sub-grid edge samples $\alpha$ on its two
sides, sees a jump, and treats the coefficient as smooth between its nodes,
which is the first-order failure of the jump case. It recovers once the grid
resolves the edge. In the one-dimensional ramp problem of
Section~\ref{sec:knee1d}, the error of the standard fourth-order operator
through a $9 : 1$ edge of width $\delta = 0.01$ or $0.0025$ falls by 100 to
200 times between $h = 2\delta$ and $h = \delta/2$, where a first-order line
falls by four. On scattered nodes in two dimensions
(Section~\ref{sec:naive2d}) its elliptic error stays at $0.86$ to $0.94$ of
its error at the jump while $h \gtrsim 6\delta$, falls to $0.081$ to
$0.084$ of it at $h \approx \delta$ for the three widths that reach it, and
converges at the fifth order of the bulk stencils below. We call the turn at
$h \approx \delta$ the knee.

% TRACE: c = (a - b) ln(a/b)/(2ab), stiff 1.7 and 2.2 (E3.3's closed form,
% heat1d.exact.edge_resistance_deficit); 0.11 at 1600 nodes and
% delta = 0.01, stiff 2.5 (1) (the two-constant medium's equilibrium);
% "only while h >~ 3 delta, by at most a factor of 6", stiff 5.3 (3).
The interface-aware stencils do not reach this regime either. Applied with
the centre of the edge as the interface and the two sides' coefficients as
if $\delta$ were $0$, they converge to the jump problem's solution, which is
not the smooth problem's: the resistance $\int dx / \alpha$ of a $\tanh$ edge
between conductivities $a$ and $b$ falls short of the jump's by $c\,\delta$,
$c = (a - b) \ln(a/b) / (2ab)$, so their error sits on a floor of order
$\delta$ while the grid does not resolve the edge. Once it does, the rebuilt
rows keep enforcing a kink that the solution lacks, and their error grows:
through a $9 : 1$ edge with $\delta = 0.01$ it reaches $0.11$ at 1600
nodes in the one-dimensional equilibrium of
Section~\ref{sec:construction1d}. On scattered nodes they
beat the standard stencils only while $h \gtrsim 3\delta$, and by at most a
factor of $6$. Using them means knowing $\delta$ and switching them off.

% TRACE: the chain, stiff 1.2 (eq. 1.2's definition; the heat polynomials,
% the even and odd seeds as the continuity conditions); 1.3 (the
% first-order form); 1.4 (the jump as the limit); the companion,
% LITERATURE.md 1b O1's last obligation (stiff 1.2: its u-field chain with
% rho = 1, K = alpha).
This paper keeps everything about the scheme except the basis of the
stencils that see the edge. In place of the monomials a stencil is made
exact on \emph{seeds}, the $t = 0$ profiles of solutions of \eqref{eq:heat}
that are polynomial in time. With a constant $\alpha$ these profiles are
the polynomials, the initial values of the heat polynomials
\cite{RosenbloomWidder1959}. In one dimension, with $L = \px \alpha \px$
and anchored at a stencil's evaluation point $x_e$, they form the chain
\begin{equation}
  \label{eq:chain}
  \phi_0 = 1, \qquad
  \alpha \phi_1' = \alpha(x_e), \qquad
  L \phi_k = k (k - 1)\, \alpha(x_e)\, \phi_{k-2} \quad (k \ge 2),
\end{equation}
with $\phi_k(x_e) = 0$ for $k \ge 1$ and $\phi_k'(x_e) = 0$ for $k \ge 2$,
integrated through the edge as a first-order system of ordinary
differential equations that never differentiates $\alpha$
(Section~\ref{sec:chain}). The seed $\phi_1$ has a constant flux, and
$\phi_2$ is the profile of a solution whose rate of change $u_t$ is
constant; in general the even seeds carry
$\partial_t^k u = \text{const}$ and the odd ones
$\alpha\, \px \partial_t^k u = \text{const}$, which are the interface
conditions of a jump. Read across a jump, the chain is the translated
polynomial basis of the interface-aware stencils
(Section~\ref{sec:jumplimit}), so one construction covers the jump, the
sub-grid edge and the resolved edge. In one dimension it is classical
(Section~\ref{sec:priorwork}). It is also the construction of our companion
paper for the wave equation \cite{Martin2026seeds}, specialised to a single
field: the two interface conditions of diffusion, on $u$ and on its flux,
are both conditions on $u$, so one chain serves every stencil.

% TRACE: stiff 3.2-3.4 (the straight feature: seeds polynomial in x' with
% coefficient functions marched in y'), 3.5 and 4.6 (a frozen normal profile
% O(1)-inconsistent where alpha varies along an unresolved edge; route (a)),
% 3.10 (the tangential chain in the curve's own coordinates), 3.11 (the
% ring's exact width, series and flux seeds); the flux-seed sentence is
% LITERATURE.md 6b's, required wherever the flux seeds are named (K9).
In two dimensions the seeds replace the polynomial augmentation of the
interface-aware RBF-FD stencils, and the Gaussian part and the
saddle-point solve stay as they were. Where $\alpha$ varies along the
normal only, as across a straight edge, seeds polynomial in the tangential
coordinate turn the chain into a triangular system of ordinary differential
equations in the normal coordinate, one short march per stencil
(Section~\ref{sec:straight}). Where the edge is curved or $\alpha$ varies
along it, a profile frozen along the normal is inconsistent at $O(1)$
through an unresolved edge, and the chain is marched in the curve's own
coordinates instead (Section~\ref{sec:curved}). A resistive layer thinner
than the stencil, the insulating ring of \cite{MartinFornberg2017heat},
needs its exact width, its own series and seeds of one degree higher that
carry the flux (Section~\ref{sec:ring}). An imperfect contact costs
strong-form schemes a derivative in the jump conditions, as Lombard and
Piraux observed for elastic waves \cite{LombardPiraux2003,LombardPiraux2006};
the flux seeds add the functions of one degree higher that carry the flux,
and nothing else.

% TRACE: the comparators, stiff 1.1 (baselines), 2.4 and 4.9 (the
% treatments), plan 3.4; the 1-D elliptic degeneracy on constant pieces,
% stiff 1.8 and 2.5 (4); the 2-D elliptic problem as a real test, stiff
% 5.3 (5).
We measure the seeds in one and two dimensions, on elliptic and parabolic
problems, against three baselines: the standard stencils with $\alpha$
sampled at the nodes; the interface-aware stencils of the jump applied at
the edge's centre; and the coefficient treatments of diffusion practice,
which change the coefficient and keep the scheme: harmonic and arithmetic
means of $\alpha$ over a cell (over a disc on scattered nodes), the edge
widened to the grid, and, in one dimension, the conservative scheme with
exact face conductances. On constant pieces the one-dimensional equilibrium
cannot rank these methods, since its solution has a constant flux and lies
in the span of the first two seeds, on which the seeds and the exact face
conductances are both exact (Section~\ref{sec:elliptic1d}); the
one-dimensional ranking is the parabolic problem's. In two dimensions the
tangential variation of the solution takes it out of that span, and the
equilibrium is a real test.

\subsection*{Contributions}
% TRACE, item by item. Item 1: LITERATURE.md 6b's contributions sentence,
% verbatim; stiff 3.2-3.4, 3.10, 3.11. Item 2: stiff 2.5 (1) (100 to 200
% and 800 to 1000 times between h = 2 delta and delta/2, checked as
% bounds; the construction on the c delta floor for h >~ 2 delta, 0.93-1.00
% of it on the two-constant medium, a check E5.4 added), 2.5 (2)
% (0.1-0.3 %, 1-1.5 %, 3-5 % off the delta = 0 line over 50-400 nodes; the
% spatial error alone fourth order, E4.10's roundtable scratch, stiff 5.6).
% Item 3: stiff 5.3 (4) (fits 4.23, 4.31, 4.23, 4.48; the spread 1.2-2.4;
% 7.99e-9 against 8.36e-5 at 40,000 nodes and delta = 0.0025; the rule's
% 20 delta, a parameter; "no deterioration ... not a proven constant", the
% statement's own caveat), 5.3 (6) (least diagonal dominance never below the
% jump-aware rows'; gmres counts below the construction's; no eigenvalue
% in the right half-plane). Item 4: stiff 5.3 (8) (1.4-3.9 from 5000
% nodes, RMS), 5.3 (9) (1.3-2.1e-15 at every s; fits 4.43-4.71 over
% 1250-40,000). Item 5: LITERATURE.md 6b's required caveat, verbatim (801;
% 1.75; 2.3 to 4.7 orders over 1250 to 40,000 nodes: stiff 2.5 (3), 5.3
% (10)). Item 6: the repository (public, MIT; paper/ CC BY 4.0); the
% version tag is #52's to pin.
\begin{enumerate}
\item The seeds for the diffusion operator on the unchanged scattered nodes
  of the interface-aware RBF-FD stencils, which they reduce to as
  $\delta \to 0$: the one-dimensional chain (Section~\ref{sec:seeds1d}), its
  two-dimensional form for a straight edge in the normal coordinate, for a
  curved or tangentially varying edge in the curve's own coordinates, and
  for a resistive layer thinner than the stencil
  (Section~\ref{sec:seeds2d}). To our knowledge no earlier scheme builds
  the basis of stencils on an unchanged scattered node set through a smooth
  material edge thinner than the node spacing by marching ordinary
  differential equations; its ingredients (the exact face conductance, the
  formal powers, bases from local solutions of the equation) are each
  classical, and in one dimension so is the construction.
\item In one dimension (Section~\ref{sec:results1d}): the knee of the
  standard operator, a drop of 100 to 200 times between $h = 2\delta$ and
  $h = \delta/2$ at a $9 : 1$ edge and of 800 to 1000 times at the
  $10 : 1$ edges of the sinusoidal medium; the jump construction on its
  floor, the distance of order $\delta$ between the jump's solution and the
  smooth one, while $h \gtrsim 2\delta$ on the two-constant medium, and
  growing once the grid resolves the edge; and the seeds, fourth order at every $\delta$ on one line: on the
  ramp problem of the two-constant medium they lie within $0.1$ to $0.3\%$,
  $1$ to $1.5\%$ and $3$ to $5\%$ of the jump's line at
  $\delta = 0.0025$, $0.01$ and $0.04$.
\item On scattered nodes through a flat band (Section~\ref{sec:results2d}):
  the seeds are fourth order at every $\delta$ tested, elliptic and
  parabolic, with elliptic fits of $4.23$ to $4.48$ over 1250 to 40,000
  nodes; the elliptic errors at the five widths, the jump included, differ
  by factors of $1.2$ to $2.4$ at each node count, and at 40,000 nodes and
  $\delta = 0.0025$ the seeds' error is $\sci{7.99}{-9}$ against the
  standard operator's $\sci{8.36}{-5}$. That is no deterioration as
  $\delta \to 0$ on these node sets, not a proven $\delta$-uniform
  constant. The rule is to seed every row whose stencil sees an edge
  within $20\delta$, with no threshold in $\delta$. The seeded rows'
  least diagonal dominance never falls below the jump-aware rows', the
  seeded matrices' condition estimates and GMRES counts stay bounded as
  $\delta \to 0$, and their spectra have no eigenvalue in the right
  half-plane.
\item Through the curved band of case 2 the tangential chain is fourth
  order, within $1.4$ to $3.9$ times of the flat seeds' RMS error from 5000
  nodes at every $\delta$ (Section~\ref{sec:curved2d}). On the insulating
  ring thinned at a fixed resistance from $s = 10^3$ to $10^{11}$ the seeds
  are exact on the matched radial profile to rounding ($\sci{1.3}{-15}$ to
  $\sci{2.1}{-15}$) and converge at fourth order (fits $4.43$ to $4.71$) at
  every $s$; through a smooth ring whose width and edges are both below $h$
  they converge at fourth order on the nodes away from it
  (Section~\ref{sec:ring2d}).
\item The comparison with the coefficient treatments
  (Sections~\ref{sec:treatments1d} and~\ref{sec:treatments2d}). The
  comparison is against coefficient sampling on the same nodes and against
  our implementation of the cell harmonic and arithmetic means, a widened
  edge and, in one dimension, the conservative scheme with exact face
  conductances, on one contrast and one test problem per dimension. Through
  an edge the grid does not resolve none of them is better than second
  order; the conservative scheme is exact at the one-dimensional
  equilibrium and, on the sinusoidal medium at $\delta = 0$, ahead of the
  seeds below 801 nodes, where the seeds' line is pre-asymptotic; on the
  scattered nodes none improves on sampling by more than a factor of 1.75
  in the RMS norm, and the seeds lie 2.3 to 4.7 orders below the best of
  them over 1250 to 40,000 nodes.
\item An open implementation in Python, written with the AI tools disclosed
  at the end of the paper, with a test for every numerical routine, in
  which every figure and table of this paper is drawn by a committed script
  from committed results files
  (\url{https://github.com/bradleypmartin/rbf-elliptic-parabolic-interfaces-2026}).
\end{enumerate}

\subsection{Relation to prior work}
\label{sec:priorwork}
% Placed from LITERATURE.md 6a (#43), cut as 6a allows, in its order: (1) the
% HODIE and fitted-operator sentence is dropped (Section 3 cites them);
% (2) the jump methods are cut to the IIM and CutFEM keys (K8: LeVeque-Li,
% Li-Ito; Hansbo-Hansbo, Burman et al.) and the two post-2016 ones; (3) the
% resistive-layer sentence moves to Section 5.4 (sec:ring, where it waits
% as a comment for #48), leaving one clause here. The first sentence, the
% FLAME and LOD clauses and the last four sentences are verbatim: they carry
% the obligations of LITERATURE.md 1a K2-K5, K10 and 1b O1, and it stays one
% paragraph because O1 asks for K3-K6 in the claim's paragraph. About 480
% words: 6a's 250-350 target is retired (its note, Brad, 2026-09-24). Any
% change to the novelty language re-opens the ledger first.
% TRACE: 1.75, stiff 5.3 (10), paper/data/heat2d_stiff_treatments*.json
% (paper_numbers.py, where "abstract, §1; stiff §5.3 (10)").
In one dimension the construction is classical in every part, and we claim
nothing for it. The first seed, the constant-flux function
$\alpha_e \int d\xi/\alpha$, is over one cell the exact face conductance of
Tikhonov and Samarskii's homogeneous schemes
\cite{TikhonovSamarskii1962,Samarskii2001} and, for a piecewise-constant
coefficient, the harmonic-mean interface conductivity of finite-volume heat
conduction \cite{Patankar1980}; the conservative scheme built on it, exact
at the source-free equilibrium, is second order \cite{PanXuLi2026}. The
whole chain is the representation of Sturm--Liouville solutions by
recursive integrals, as power series in the spectral parameter
\cite{Kravchenko2008,KravchenkoPorter2010}, taken from the anchor; for a
constant coefficient its time-polynomial solutions are the heat polynomials
\cite{RosenbloomWidder1959,Widder1975}. Three-point schemes of any order
for $(k u')'$ with a piecewise-smooth $k$, whose coefficients come from
Cauchy problems integrated across each cell, are the exact and truncated
difference schemes \cite{SamarskiiMakarov1990,GavrilyukEtAl2011};
three-point schemes of any order exist across a jump too
\cite{Gartland1989}, one of them built by carrying Taylor data across the
step \cite{Sujecki2010}. In two dimensions, local solutions of the equation that satisfy the
interface conditions are the basis of the Trefftz difference schemes of
FLAME, on irregular stencils too \cite{Tsukerman2006,Tsukerman2010}, and of
the tailored finite point methods \cite{Huang2009,HanZhang2012}; Galerkin
methods build their bases from local solves
\cite{BabuskaCalozOsborn1994,HouWu1997,ChuGrahamHou2010,OwhadiZhang2007,MalqvistPeterseim2014}
and reach high order through coefficients the mesh does not resolve
\cite{Maier2021,HauckPeterseim2023}, and the quasi-Trefftz spaces, one of
them for the heat equation, build polynomial local solutions by Taylor
expansion \cite{ImbertGerardMoiolaPerinatiStocker2025,GomezPerinatiStocker2026},
which cannot reach through an edge thinner than the stencil
\cite{Martin2026seeds}. For a jump, the interface-aware RBF-FD stencils
this work generalises \cite{Martin2016,MartinFornberg2017heat} stand beside
the immersed interface method and unfitted finite elements
\cite{LeVequeLi1994,LiIto2006,HansboHansbo2002,BurmanEtAl2015}, now of
fourth and sixth order \cite{FengZhao2020,FengHanMinev2024}, all of which
take the jump as given. A resistive layer thinner than the grid is an
imperfect interface in the limit
\cite{SanchezPalencia1974,PerrusselPoignard2013,Hashin2001}, which
Section~\ref{sec:ring} takes up. The standing alternative changes the
coefficient rather than the stencil: harmonic or arithmetic means over a
cell \cite{Patankar1980,EwingIlievLazarov2001,GoodeAppel1992}, averaging for
layering finer than the grid \cite{MoskowEtAl1999,DavydychevaDruskinHabashy2003},
or a regularised coefficient \cite{TornbergEngquist2003,TornbergEngquist2006}.
Those methods take the cell or the width from the grid; here the width is a
property of the medium, below the node spacing, and the aim is the scheme's
full order through it. We found no prior instance of stencils on an
unchanged scattered node set whose basis is continued through a sub-grid
smooth edge by ordinary differential equations, nor of such stencils
reducing to a jump construction as $\delta \to 0$; the search is logged in
the repository's \texttt{LITERATURE.md}. We ran the cell means, a widened
edge and, in one dimension, the conservative scheme with exact face
conductances, as our implementation of each on one problem per dimension:
through an edge the grid does not resolve none is better than second order,
and on the scattered nodes none improves on sampling the coefficient by
more than a factor of $1.75$ in the root-mean-square error. We claim nothing
beyond that measurement.

\subsection*{Outline}
Section~\ref{sec:setting} fixes the problems, the discretisations and the
interface-aware stencils of a jump that the seeds extend.
Section~\ref{sec:seeds1d} constructs the seeds in one dimension and proves
what can be proved about them. Section~\ref{sec:results1d} measures them
against the standard operator, the jump construction and the coefficient
treatments. Section~\ref{sec:seeds2d} builds the seeds on scattered nodes
for straight, curved and doubly sub-grid features.
Section~\ref{sec:results2d} measures them there, elliptic and parabolic,
with the matrices they make. Section~\ref{sec:discussion} records the
limitations and what we would do next.

% ===========================================================================
\section{Setting}
\label{sec:setting}
% TRACE: stiff 1.1 (1-D) and 3.1 (2-D); the dissertation ch. 4-5 and EABE
% 2017 sections 2-3 for the port, quoted, not re-derived (page map in
% docs/paper-index.md); docs/port-notes.md sections 1-2 for what the port
% reproduces. Drafted by #45. Numbers are traced paragraph by paragraph.
%
% NOTATION fixed here for the later sections (#46-#50): n nodes and h the
% spacing in both dimensions; delta the edge width; sigma the blend weight
% (s is EABE's ring parameter); alpha^- and alpha^+ an edge's pieces; L the
% operator; x_e the anchor and (x', y') the local frame; phi_k the 1-D
% seeds; the two 1-D media by name ("two-constant", "sinusoidal") and the
% ramp problem; the operators by name (naive, direct, jump-aware, the
% delta = 0 construction, seeds). The complete notation list is #51's.

\para{Notation}
The symbols the sections share, fixed here:
\begin{center}
\small
\begin{tabular}{@{}lp{0.64\textwidth}@{}}
\toprule
$n$, $h$ & node count and spacing: $h = 2 / (n - 1)$ in one dimension and
  the row spacing $1 / \operatorname{round}(0.95 \sqrt{n})$ in two \\
$\alpha$; $\alpha^-$, $\alpha^+$; $L$ & conductivity; the pieces on the two
  sides of an edge; $L = \nabla \cdot (\alpha \nabla\, \cdot)$ \\
$\delta$; $\sigma$ & width of a $\tanh$ edge, $\delta = 0$ the jump; the
  blend weight $\sigma(z) = \tfrac12 (1 + \tanh z)$ \\
$s$ & the thinness of the ring of \cite{MartinFornberg2017heat}, $1/s$ its
  width \\
$x_e$; $(x', y')$ & a stencil's evaluation point (anchor); its local frame,
  $y'$ normal to the edge \\
$\phi_k$ & the one-dimensional seeds \eqref{eq:chain} \\
$D_x$, $D_y$; $A$ & first-derivative matrices; $\operatorname{diag}
  \alpha$ at the nodes \\
naive, direct & $D_x A D_x$ ($+ D_y A D_y$); $\alpha \nabla^2 u + \nabla
  \alpha \cdot \nabla u$ on one stencil \\
jump-aware & the interface-aware stencils of Section~\ref{sec:jump} \\
the $\delta = 0$ construction & the same, applied at a smooth edge's
  centre \\
seeds & the operator whose rows that see an edge are rebuilt on seeds \\
\bottomrule
\end{tabular}
\end{center}

\subsection{The problem and its interface conditions}
\label{sec:problem}
% TRACE: stiff 1.1 (eq. 51-55 of the dissertation, EABE eq. 1: continuity
% of D^k u and alpha d_x D^k u; the tanh edge; the three regimes), 3.1
% (the 2-D edge in the signed normal distance, composed edge by edge);
% 2.1 (the tails: within two ulps of the pieces at 19 delta, the pieces
% to the bit beyond 20 delta); c = (a - b) ln(a/b)/(2ab), stiff 1.7 and
% 2.2; 8.79 for the 9 : 1 edge, stiff 2.5 (1), paper/data/heat1d_stiff.json
% (paper_numbers.py, where "§2.1; stiff §2.5 (1)"); the norms and fits,
% stiff 5.3's preamble and 5.5, port notes 2.10.
On a domain $\Omega$ in one or two dimensions we solve \eqref{eq:heat} with
Dirichlet data on $\partial \Omega$, and its equilibrium $L u = 0$, the
elliptic problem. The conductivity $\alpha > 0$ is one smooth function on
each side of every interface $\Gamma$. Across a jump the temperature and the
heat flux are continuous at all times, hence so are all their time
derivatives, and $\partial_t^k u = L^k u$ turns these into conditions on
spatial derivatives,
\begin{equation}
  \label{eq:continuity}
  [L^k u] = 0, \qquad [\alpha\, \partial_n L^k u] = 0 \qquad \text{on }
  \Gamma, \quad k = 0, 1, 2, \ldots,
\end{equation}
where $[\cdot]$ is the difference of the two sides' limits, $n$ the normal
to $\Gamma$, and each side's $L$ is taken with that side's $\alpha$
\cite{Martin2016,MartinFornberg2017heat}.

\para{A smooth edge}
With an edge width $\delta > 0$ the jump at an interface is replaced by a
$\tanh$ transition between the two pieces,
\begin{equation}
  \label{eq:edge}
  \alpha = \bigl(1 - \sigma(d / \delta)\bigr)\, \alpha^- + \sigma(d / \delta)\,
  \alpha^+, \qquad \sigma(z) = \tfrac12 (1 + \tanh z),
\end{equation}
where $d$ is the signed distance to the interface, $x - x_c$ in one
dimension and the distance along the normal through the foot point on the
curve in two, positive on the side of $\alpha^+$. The pieces keep their own
variation, and $\delta = 0$ is the jump. A band between two interfaces is
blended twice, from the outside piece into the inside one across the lower
interface and back across the upper, so that a band wider than its edges
has two independent transitions; Section~\ref{sec:ring} defines the smooth
form of a layer thinner than its edges. Beyond $19\delta$ from its centre
an edge is its pieces to within two units in the last place.

\para{Three regimes}
For $\delta \gg h$ a stencil sees a smooth coefficient. For $\delta = 0$ it
sees a jump, and \eqref{eq:continuity} applies. For $0 < \delta \lesssim h$
it sees what looks to it like a jump between two of its nodes, but no
interface condition applies, since $\alpha$ is smooth, and the smooth
problem and the jump problem have different solutions. They differ at
first order in $\delta$: in one dimension the resistance $\int dx / \alpha$
of a $\tanh$ edge between constant pieces $a$ and $b$ falls short of the
jump's by
\begin{equation}
  \label{eq:deficit}
  c\, \delta, \qquad c = \frac{(a - b) \ln (a / b)}{2 a b},
\end{equation}
which is $8.79\,\delta$ for the $9 : 1$ edge of Section~\ref{sec:setting1d}
and is positive for either ordering of $a$ and $b$.

\para{Norms and orders}
In one dimension an error is the relative $\ell^2$ error at the nodes,
$\|e\|_2 / \|u\|_2$, and a rate is $\log_2$ of the ratio of consecutive
errors as the node count doubles. In two dimensions an error is the
root-mean-square (RMS) error over all nodes, the Dirichlet nodes included,
and an order is the least-squares slope of $\log e$ against $\log h$ over
the node counts named. A maximum-norm error is named as such wherever one
is quoted, and every caption says which norm it shows.

\subsection{One dimension}
\label{sec:setting1d}
% TRACE: the grid, port notes 1.1 (h = 2/(n - 1), both ends on nodes;
% eq. 75's edges on nodes when (n - 1) % 4 == 0); the media, stiff 1.1 and
% 2.1 (the two-constant medium 1/9 | 1 of heatEq1DMatlab/FD4heat1DAC.m, the
% edge mid-cell at even n; eq. 75's layer with both edges at 1 | 0.1); the
% four widths, stiff 2.2 (STUDY_DELTAS); the ramp problem, port notes 1.4
% and 1.6 (the C-infinity step over [0, 1], read at t = 2; the MATLAB's
% logistic ramp incompatible at the corner), scripts/heat1d_stiff.py (RAMP,
% T_END, BC); the naive operator, port notes 1.1 (nine-point rows, one-sided
% FD4 at the ends); the direct stencil at 0.22 on eq. 75, port notes 1.3;
% BD4, port notes 1.6 (RK4 start-up at half the infinity-norm limit,
% STARTUP_FRACTION; Dirichlet values at every stage) and 2.5 (the 2-D
% analytic history).
The one-dimensional problem is \eqref{eq:heat} with $L = \px \alpha \px$ on
$[-1, 1]$ with Dirichlet ends. A grid has $n$ equispaced nodes with both
ends on nodes, so $h = 2 / (n - 1)$. Two media are used, each at
$\delta = 0$, $0.04$, $0.01$ and $0.0025$:
\begin{itemize}
\item the \emph{two-constant medium}, $\alpha = 1/9$ for $x < 0$ and $1$ for
  $x > 0$, the single interface of the time-dependent test in the code
  behind \cite{Martin2016}, which the dissertation does not report; the
  edge lies midway between two nodes when $n$ is even;
\item the \emph{sinusoidal medium} of \cite[eq.~(75)]{Martin2016},
  $\alpha = 0.1 + 0.4 \sin 2\pi x$ on the layer $[0, \tfrac12]$ and $1$
  elsewhere, whose two edges both have the contrast $10 : 1$, since the
  sinusoid vanishes at both; both edges fall on nodes when $n - 1$ is a
  multiple of $4$.
\end{itemize}
Each medium carries two problems: the equilibrium with $u(-1) = 1$ and
$u(1) = 0$, and the \emph{ramp problem}, which starts from $u = 0$, holds
$u(1, t) = 0$, raises $u(-1, t)$ from $0$ to $1$ over $0 \le t \le 1$ by
the $C^\infty$ step
$e^{-1/t} / (e^{-1/t} + e^{-1/(1-t)})$, and is read at $t = 2$. (The
original code ramps with a logistic function, whose value and derivatives at
$t = 0$ are small but not zero; the mismatch with the zero start at the
corner would cap a fourth-order study.) Section~\ref{sec:media1d} gives the
reference solutions.

The \emph{naive} operator is $D_x A D_x$ \cite[eq.~(76)]{Martin2016}, with
$D_x$ the fourth-order first-derivative matrix of Fornberg's weights
\cite{Fornberg1988}, one-sided at the ends, and $A = \operatorname{diag}
\alpha(x_j)$, which gives nine-point rows in the interior. It is the form in
which a standard stencil sees a jump at all: a direct discretisation of
$\alpha u_{xx} + \alpha' u_x$ with the coefficients sampled at the nodes is
blind to the jumps, and its error on the sinusoidal medium stays at $0.22$
as the grid is refined. Time stepping is the fourth-order backward
differentiation formula (BD4) with $\Delta t = h$, one sparse LU
factorisation per operator and the Dirichlet values imposed at every step;
its three starting values come from classical Runge--Kutta sub-steps at half
the stability limit.

\subsection{Two dimensions}
\label{sec:setting2d}
% TRACE: the cases, port notes 2 (EABE eq. 32-40 via docs/paper-index.md;
% case 2's curves from the MATLAB's curvedinterface1/2, the index's note)
% and stiff 3.1; case 1's exact solutions, port notes 2.2 and 2.5
% (heat2d.exact.case1_exact: v(0) = 0, v(1) = 1; e^t sin 2 pi x v(y) to
% t = 0.1); the ring's resistance 1.5 at every s and contrast about
% 1500 : 1, port notes 2.9 and stiff 3.1; the node sets, port notes 2.1
% (m = round(0.95 sqrt n); rows at 0.5, 0.5 + sqrt 3/2, 0.5 + sqrt 3 row
% spacings; 100 repulsion steps; the ring's midline layout); RBF-FD, port
% notes 2.2 (epsilon = 0.4/d with d the nearest neighbour; 42 / 5 and
% 30 / 4; the boundary zone 3/sqrt n; the interface group by the crossing
% test, 2.3); the direct operator blind to the jump (3.6e-2 at every count
% on case 1, not quoted), port notes 2.2; solvers, port notes 2.8 (the
% reduced interior system; spilu; Appendix B's P); BD4 from the analytic
% history, 3 to 19 steps to t = 0.1, port notes 2.5 and stiff 5.5.
The two-dimensional problems are the three test cases of
\cite{MartinFornberg2017heat}, \eqref{eq:heat} on the unit strip $0 \le y \le 1$, periodic in $x$ with
period $1$, with Dirichlet data on $y = 0$ and $y = 1$. The conductivity is
one smooth piece in a band between two interfaces and another outside it.
\begin{itemize}
\item \emph{Case 1}, a flat band $0.6 \le y \le 0.8$ with $\alpha = 0.2$ in
  it and $1$ outside. With $u = \sin 2\pi x$ on $y = 1$ and $u = 0$ on
  $y = 0$ the equilibrium is $\sin 2\pi x\, v(y)$ with $v$ piecewise
  exponential, and the parabolic problem is marched to $t = 0.1$ against the
  exact solution $e^t \sin 2\pi x\, w(y)$ of the same form
  \cite{MartinFornberg2017heat}.
\item \emph{Case 2}, both interfaces bent to $y = 0.6 + 0.02 \sin 2\pi x$
  and $y = 0.8 + 0.02 \sin 2\pi x$,\footnote{The published texts do not
  state the curves; these are the ones of the original code.} with
  $\alpha = 0.2 + 0.1 \sin 2\pi x \sin 2\pi y$ in the band and the boundary
  data of case 1.
\item \emph{Case 3}, a ring $0.35 - 1/s \le r \le 0.35$ about
  $(\tfrac12, \tfrac12)$ with $\alpha = 1/(1.5 s) + \sin 2\pi x \sin 2\pi y
  / (3 s)$ in it and $1$ outside, around a cooling circle $r = 0.05$ whose
  inside is removed, with $u = \sin 6\pi x$ on $y = 0$ and $y = 1$ and
  $u = 0$ on the circle. At $s = 10^3$ this is the insulating ring of
  \cite{MartinFornberg2017heat}, $0.001$ wide at a contrast of about
  $1500 : 1$; the extremizing sweep of that paper runs $s$ up to $10^{11}$,
  thinning the ring at a fixed resistance, its width over its conductivity,
  of $1.5 / (1 + \tfrac12 \sin 2\pi x \sin 2\pi y)$.
\end{itemize}
With $\delta > 0$ each interface of cases 1 and 2 becomes an edge
\eqref{eq:edge} in the signed normal distance to its curve; the ring's
smooth form is Section~\ref{sec:ring}'s. On case 1 the medium is then the
one-dimensional one in $y$, and the reference solutions stay separable
(Section~\ref{sec:problems2d}).

\para{Node sets}
A set of $n$ nodes is built as in \cite{MartinFornberg2017heat}. Rows hold
$m = \operatorname{round}(0.95 \sqrt n)$ nodes per unit length, so their
spacing is $h = 1/m$. Six fixed rows straddle each interface, three on each
side at normal distances $(\tfrac12 + j \tfrac{\sqrt 3}{2}) h$,
$j = 0, 1, 2$, the middle row on each side staggered by half a spacing
along the curve, so that the rows cross the interface in a hexagonal
pattern; a fixed row
lies on each Dirichlet curve; and the remaining nodes are placed by 100
steps of a simulated repulsion from a random start, kept clear of the
straddling rows. The ring is thinner than $h$ at every node count used, so
one set of six rows straddles its midline, and both of its interfaces pass
between the innermost pair. Unless a spread over node sets is quoted,
each node count has one node set, from a fixed random seed.

\para{RBF-FD}
RBF-FD \cite{FornbergFlyer2015primer,FornbergFlyer2015acta} builds each row
of a differentiation matrix from a stencil. For a stencil of nodes
$x_1, \ldots, x_{N_s}$ about an evaluation node $x_c$, the weights $w$ of a
linear operator $\mathcal{L}$ solve
\begin{equation}
  \label{eq:rbffd}
  \begin{bmatrix} \Phi & P \\ P^{\mathsf T} & 0 \end{bmatrix}
  \begin{bmatrix} w \\ w_* \end{bmatrix}
  =
  \begin{bmatrix}
    \mathcal{L}\, \varphi(\|x - x_i\|) \big|_{x_c} \\
    \mathcal{L}\, p_j(x) \big|_{x_c}
  \end{bmatrix},
\end{equation}
with $\Phi_{il} = \varphi(\|x_i - x_l\|)$ for the Gaussian
$\varphi(r) = e^{-(\varepsilon r)^2}$, $P_{ij} = p_j(x_i)$ over the
bivariate monomials $p_j$ up to a degree $p$ in coordinates scaled by the
stencil's radius, and $w_*$ discarded. The constraint rows make the weights
exact on the polynomials, which take over from the Gaussians under
refinement \cite{FlyerFornbergBayonaBarnett2016,BayonaFlyerFornbergBarnett2017};
the shape parameter is $\varepsilon = 0.4 / d$, with $d$ the distance from
$x_c$ to its nearest neighbour. Stencils have $42$ nodes and degree $5$ in
the bulk, and $30$ nodes and degree $4$ within $3 / \sqrt n$ of a Dirichlet
curve and wherever a $42$-node stencil would reach across an interface.
Of the latter, the interface-aware construction rebuilds the rows whose
$30$ nodes cross an interface (Section~\ref{sec:jump}); the seeds rebuild
the rows whose stencils see an edge (Section~\ref{sec:assembly2d}).

\para{Operators and solvers}
The naive operator is $D_x A D_x + D_y A D_y$, the two-dimensional form of
\cite[eq.~(76)]{Martin2016}, with every node's rows of $D_x$ and $D_y$ built
on its own stencil, so that it reaches the neighbours of the neighbours.
The direct operator applies $\alpha \nabla^2 u + \nabla \alpha \cdot \nabla
u$ at each node through one stencil; it is blind to a jump and serves as
the control where the grid resolves the edge. Dirichlet rows are the
identity. The elliptic systems are solved directly with SuperLU; the
iterative study applies GMRES and BiCGSTAB to the interior system, with and
without an incomplete LU factorisation and the diagonal-dominance
preconditioner of \cite[App.~B]{Martin2016}. The parabolic problem of
case 1 is marched by BD4 with $\Delta t = h$ from the exact solution's
history, $3$ to $19$ steps to $t = 0.1$.

\subsection{Interface-aware stencils for a jump}
\label{sec:jump}
% TRACE: port notes 1.1 (the 1-D translated basis: degree 4, EABE eq. 24,
% three or four rebuilt rows per interface; two interfaces translate
% twice, 1.6), 2.3 (the 2-D continuity rows: the first p + 1 - 2k
% coefficients of D^k u and p - 2k of the flux, 15 for p = 4; the
% curvature entering twice, EABE eq. 29-30 / dissertation eq. 82-83; the
% frame at the nearest point, in stencil-radius units), 2.4 (the warp:
% slope ratio alpha^- / alpha^+, value and alpha d_n continuous, first
% order in the RBF part; EABE 2.2.4); stiff 1.2 (the even and odd seeds as
% the five conditions), 1.4 and 5.3 (4) (the jump limit); stiff 1.1 and 3.1
% (the delta = 0 construction as named).
The seeds extend the interface-aware stencils of
\cite{Martin2016,MartinFornberg2017heat}, whose wave-equation form is in
\cite{MartinFornbergStCyr2015,MartinFornberg2017}. We summarise them in
the form the later sections use; the derivations are in those references.

\para{One dimension}
Expand $u$ on each side of an interface $x_c$ in a polynomial of degree
four about $x_c$, with the Taylor coefficients of that side's $\alpha$
entering through multiplication matrices on the coefficient vectors, and
impose \eqref{eq:continuity} on $L^k u$ for $k = 0, 1, 2$ and on the flux
for $k = 0, 1$. These five conditions give a square \emph{continuity
matrix} per side, $C^- a^- = C^+ a^+$ on the coefficient vectors $a^\pm$.
Keeping the monomials on the evaluation node's side and carrying them
across with $(C^+)^{-1} C^-$ gives the \emph{translated basis}: piecewise
polynomials, one per monomial, that satisfy the five conditions. The weights
of a five-node stencil that straddles $x_c$ reproduce $L$ at the evaluation
node on this basis, a $5 \times 5$ solve per stencil, and a stencil that
straddles both edges of a thin layer translates twice.

\para{Two dimensions}
A stencil whose $30$ nodes lie in more than one region is given a local
frame at the point of each interface nearest its centre, $x'$ along the
tangent and $y'$ along the normal, in units of the stencil's radius, and
the interface is expanded there as $y' = f(x')$. The continuity conditions
equate, along the interface, the first $5 - 2k$ coefficients in $x'$ of
$L^k u$ and the first $4 - 2k$ of the normal flux
$\alpha (\partial_{y'} - f' \partial_{x'}) L^k u$, fifteen conditions for
degree four, one per bivariate monomial; the curvature enters through $f$,
in the curve along which the coefficients are read and in the normal
derivative. The translated basis replaces the monomials $p_j$ in
\eqref{eq:rbffd}. The Gaussians are \emph{warped}: in the frame of the
evaluation node's side, the normal coordinate of every node is replaced by
a piecewise-linear one whose slope changes by the factor $\alpha^- /
\alpha^+$ across each interface, so that a Gaussian's value and its
$\alpha\, \partial_n$ are both continuous there and the RBF part of the
stencil satisfies the flux condition to first order.

This construction, with warped Gaussians, is what we call
\emph{jump-aware}. Applied with the centre of a smooth edge as the
interface and the pieces' Taylor data as if $\delta$ were $0$, it is the
\emph{$\delta = 0$ construction} of Sections~\ref{sec:construction1d}
and~\ref{sec:naive2d}. It is also the $\delta \to 0$ limit of the seeds,
which reproduce it at $\delta = 0$ (Sections~\ref{sec:jumplimit}
and~\ref{sec:assembly2d}).

\subsection{Scope of the implementation}
\label{sec:scope}
% TRACE: port notes 2.10 (the reproduction table: every 2016 line at its
% 2016 rate or better; the six discrepancies, of which 4 and 5 are quoted:
% bicgstab converges on case 3 at every count, and every s of eq. 40 is
% fourth order to its reference's floor where Fig. 19 floors at 7.4e-4 for
% s = 1e11; the 2016 2-D code not preserved), 1.2-1.5 and 2.1-2.9 (the
% figures reproduced); E2.10 (#24) closed won't-do (the corner, dissertation
% 5.4.5, second order by circular segments); stiff 5.5 (what is not
% covered: corners, three dimensions, edges meeting a Dirichlet boundary,
% time-dependent media; one tanh profile; one contrast per case; 30 / 4
% seeded rows in a 42 / 5 bulk); contrasts: 9 : 1 and 10 : 1 (stiff 1.1),
% 5 : 1 (stiff 3.1), about 1500 : 1 on the ring (stiff 3.1).
The experiments run on a reimplementation of the 2016--2017 methods in
Python (NumPy and SciPy, with a test for every numerical routine), written
for this study. Every published figure of the one- and two-dimensional
heat problems of \cite{Martin2016,MartinFornberg2017heat} but the cornered
interface's has a counterpart in it, and every published convergence line
is reproduced at its published rate or better. Where the numbers part
company the repository's port notes say so. Two differences
bear on this paper: the iterative solvers do not break down on the ring
here (BiCGSTAB converges at every node count), and the extremizing sweep
does not lose convergence at $s = 10^{11}$ (every $s$ is fourth order down
to its reference's own error). The 2016 two-dimensional code is not
preserved, so neither difference could be traced to it.

The implementation leaves out the cornered interface of
\cite[Sec.~5.4.5]{Martin2016}, which is approximated there by circular
segments at second order. Its two-dimensional media are a band between two
curves, periodic in $x$, so no edge meets the domain's boundary. The bulk
stencils are degree five and the stencils that see an edge degree four,
seeded or interface-aware. Each case has one contrast: $9 : 1$ and
$10 : 1$ in one dimension, $5 : 1$ on cases 1 and 2, and on the ring about
$1500 : 1$ at $s = 10^3$, growing with $s$ at a fixed resistance. Every
smooth edge is a $\tanh$.
Section~\ref{sec:discussion} returns to these restrictions.

% ===========================================================================
\section{Seeds for the diffusion operator in one dimension}
\label{sec:seeds1d}
% TRACE: stiff 1.2 (the chain), 1.3 (its first-order form), 1.4 (nested
% kernels, the jump limit), 1.5 (the Chebyshev property), 1.6 (truncation
% error), 1.8 (the elliptic degeneracy and its finite-volume twin), with
% the measurements of 2.3 (P4-P8) and 2.5 (the seed functions); E5.5's
% notes in 1.2, 1.4, 1.5, 1.8 and 2.3; LITERATURE.md 1a K1-K3 and K7 for
% the citations (#43); code heat1d/stiff.py. Drafted by #46. Numbers are
% traced paragraph by paragraph; every cache-backed one is a check in
% scripts/paper_numbers.py whose `where` names a "§3.x".
Fix a stencil of the grid of Section~\ref{sec:setting1d}: five nodes
$x_1 < \cdots < x_5$ and an evaluation point $x_e$, one of them. The
standard weights are exact on the monomials $(x - x_e)^k$, $k \le 4$, and
the interface-aware weights of Section~\ref{sec:jump} on the translated
basis. This section builds the basis for a smooth edge, shows that the
translated basis is its limit as $\delta \to 0$, proves that every stencil
solve is nonsingular, and says what the truncation error is. Each
ingredient is classical in one dimension (Section~\ref{sec:priorwork}), and
the construction is that of our companion paper for the wave equation
\cite{Martin2026seeds}, specialised to the diffusion operator.

\subsection{Seeds as the profiles of solutions polynomial in time}
\label{sec:chain}
% TRACE: stiff 1.2 (the time-polynomial solutions, (j + 1) g_{j+1} = L g_j;
% one power of L per power of t against the wave equation's two; the heat
% polynomials of Rosenbloom-Widder; the definition; the seeds one at a
% time; the five continuity conditions in order, one field). Lemma
% lem:heatpoly's general u_k is E5.5's, added to stiff 1.2 with this
% ticket (the chain's algebra, no number).
Look for solutions of $u_t = L u$, $L = \px \alpha \px$, that are
polynomial in time, $u = \sum_{j = 0}^{J} t^j g_j(x)$. Matching powers of
$t$ gives $(j + 1)\, g_{j+1} = L g_j$ and $L g_J = 0$, so the profile at
$t = 0$ satisfies $L^{J+1} g_0 = 0$. The heat equation spends one power of
$L$ on each power of $t$ where the wave equation spends two, but the
profiles fill the same nested kernels of $L$. With a constant $\alpha$
they are the polynomials: the time-polynomial solutions of
$u_t = \alpha u_{xx}$ are the heat polynomials
\begin{equation*}
  v_k(x, t) = \sum_{0 \le j \le k/2} \frac{k!}{j!\, (k - 2j)!}\,
  x^{k - 2j}\, (\alpha t)^j
\end{equation*}
\cite{RosenbloomWidder1959,Widder1975}, whose profiles at $t = 0$ are the
monomials. The seeds are the same profiles for a variable $\alpha$,
anchored at the evaluation point.

\begin{definition}[Seeds]
\label{def:seeds}
Let $\alpha > 0$ and $\alpha_e = \alpha(x_e)$. The \emph{seeds} at $x_e$
are the functions $\phi_0, \phi_1, \ldots$ of the chain \eqref{eq:chain}:
$\phi_0 = 1$, $\alpha \phi_1' = \alpha_e$ and
$L \phi_k = k (k - 1)\, \alpha_e\, \phi_{k-2}$ for $k \ge 2$, with
$\phi_k(x_e) = 0$ for $k \ge 1$ and $\phi_k'(x_e) = 0$ for $k \ge 2$.
\end{definition}

\begin{lemma}
\label{lem:heatpoly}
For every $k$ the function
\begin{equation}
  \label{eq:heatpoly}
  u_k(x, t) = \sum_{0 \le j \le k/2} \frac{k!}{j!\, (k - 2j)!}\,
  (\alpha_e t)^j\, \phi_{k - 2j}(x)
\end{equation}
solves $u_t = L u$, and $u_k(\cdot, 0) = \phi_k$. With a constant
$\alpha$, $\phi_k = (x - x_e)^k$ and $u_k(x, t) = v_k(x - x_e, t)$.
\end{lemma}

\begin{proof}
Write $c_j = k! / (j!\, (k - 2j)!)$. By the chain,
$L u_k = \sum_j c_j\, (k - 2j)(k - 2j - 1)\, \alpha_e^{j+1}\, t^j\,
\phi_{k-2j-2}$ and $\pt u_k = \sum_j (j + 1)\, c_{j+1}\, \alpha_e^{j+1}\,
t^j\, \phi_{k-2j-2}$, and the two agree because
$(j + 1)\, c_{j+1} = c_j\, (k - 2j)(k - 2j - 1)$. With $\alpha$ constant,
$L (x - x_e)^k = \alpha\, k (k - 1) (x - x_e)^{k-2}$ and the anchor
conditions fix the two constants of integration, so $\phi_k = (x - x_e)^k$
by induction; with $\alpha_e = \alpha$, \eqref{eq:heatpoly} is then
$v_k(x - x_e, t)$.
\end{proof}

Read one at a time, the seeds say what each monomial becomes through the
edge. The seed $\phi_1$ solves $L \phi_1 = 0$: it is the equilibrium of
constant flux, $\phi_1(x) = \alpha_e \int_{x_e}^{x} d\xi / \alpha(\xi)$,
and what stays smooth through the edge is its flux $\alpha \phi_1'$, not
its slope. The seed $\phi_2$ is the profile of $u_2 = \phi_2 + 2 \alpha_e
t$, whose rate of change is constant, $u_t = 2 \alpha_e$; the constant
matters, since $u_t = 0$ would only return $\ker L$. In general, by
\eqref{eq:heatpoly}, $\pt^m u_{2m} = (2m)!\, \alpha_e^m$ is constant and
$\pt^m u_{2m+1} = (2m + 1)!\, \alpha_e^m\, \phi_1$ has a constant flux:
the even seeds carry $\pt^m u = \text{const}$ and the odd ones
$\alpha\, \px \pt^m u = \text{const}$. Taken in the order
$\phi_0, \ldots, \phi_4$ these are the five conditions of
\eqref{eq:continuity} from which the translated basis of a jump is built,
on $u$, its flux, $u_t = L u$, the flux of $u_t$, and $u_{tt} = L^2 u$. In
the wave equation the same alternation passes between two fields; in
diffusion both kinds of condition bear on $u$, so one chain serves every
stencil.

% TRACE: stiff 1.5's closing note and LITERATURE.md 1a K2 (the identity
% with q = 0, u_0 = r = 1, p = alpha, x_0 = x_e; the recursive integrals
% without a factor n, X^(k) for odd k and X~^(k) for even k, their weights
% 1/alpha and 1 alternating; the lambda-series coefficients the
% time-polynomial profiles). Checked in
% tests/heat1d/test_stiff.py::test_the_seeds_are_the_formal_powers_of_the_
% sturm_liouville_operator at 1e-9 relative through the two-constant
% medium's delta = 0.01 edge (the E5.2 pass found 12 digits for k <= 5).
% eq:nestedint is that identity with the recursion unrolled.
In one dimension the seeds are a classical object. Integrating the chain
twice from the anchor gives
$\phi_k(x) = k (k - 1)\, \alpha_e \int_{x_e}^{x} \alpha(s)^{-1}
\int_{x_e}^{s} \phi_{k-2}(r)\, dr\, ds$, and unrolling the recursion,
\begin{equation}
  \label{eq:nestedint}
  \phi_k = k!\, \alpha_e^{\lceil k/2 \rceil}\, I_k, \qquad
  I_k(x) = \int_{x_e}^{x} \frac{ds_1}{\alpha(s_1)} \int_{x_e}^{s_1} ds_2
  \int_{x_e}^{s_2} \frac{ds_3}{\alpha(s_3)} \cdots,
\end{equation}
where $I_0 = 1$ and $I_k$ has $k$ nested integrals whose weights alternate
between $1/\alpha$ and $1$, starting from $1/\alpha$ outermost. These are
the recursive integrals, $X^{(k)}$ for odd $k$ and $\tilde X^{(k)}$ for even
$k$, by which Kravchenko and Porter write the solutions of $(p u')' + q u =
\lambda r u$ as power series in the spectral parameter
\cite{Kravchenko2008,KravchenkoPorter2010}, here with $q = 0$, $r = 1$,
$p = \alpha$, the particular solution $u_0 = 1$ and the integrals taken
from the anchor. The seeds are thus the formal powers of the
Sturm--Liouville operator $L$ anchored at the evaluation point, and the
coefficients of those series are the time-polynomial profiles, with
$\lambda$ in the role of $\pt$. The repository's tests check the identity
against the ODE march below to a relative $10^{-9}$ through a $\tanh$ edge.

% TRACE: stiff 1.2 (the moment conditions, the same ones E1.2 imposes on
% the translated basis); 2.3 P4 (constant alpha = 0.7: 0.7 times
% Fornberg's weights to 1e-15 centred and 2e-16 off-centre, tests);
% LITERATURE.md 6a's HODIE and fitted-operator sentence, verbatim (#45's
% first cut dropped it from section 1 because "Section 3 cites them"); K3
% (Samarskii-Makarov, Gavrilyuk et al.: the seed stencil "a variant of
% them", K3's framing).
The weights $w_i$ of a stencil approximate $L u$ at $x_e$ from the nodal
values by reproducing $L$ on the first five seeds,
\begin{equation}
  \label{eq:moments}
  \sum_{i=1}^{5} w_i\, \phi_k(x_i) = (L \phi_k)(x_e) =
  \begin{cases} 2 \alpha_e, & k = 2, \\ 0, & k = 0, 1, 3, 4, \end{cases}
\end{equation}
since $L \phi_k = k (k - 1)\, \alpha_e\, \phi_{k-2}$ vanishes at $x_e$ for
every $k$ but $2$. These are the conditions the interface-aware stencils
impose on the translated basis, and with a constant $\alpha$ they return
$\alpha$ times the standard second-derivative weights \cite{Fornberg1988},
which the implementation reproduces to rounding. Stencils exact on a local
space other than the polynomials go back to fitted operators and to HODIE
\cite{AllenSouthwell1955,Ilin1969,LynchRice1980,RoosStynesTobiska2008}.
Three-point schemes whose coefficients come from Cauchy problems
integrated across a piecewise-smooth coefficient are the exact and
truncated schemes \cite{SamarskiiMakarov1990,GavrilyukEtAl2011}, and the
seed stencils are a five-point variant of them.

\begin{figure}
  \centering
  % figures/heat1d_stiff_seeds.pdf: scripts/paper_figures.py from
  % paper/data/heat1d_stiff.json (seed_functions, seed_functions/curves);
  % never hand-edited. TRACE: stiff 2.5 (the window: 200 nodes, h =
  % 0.01005, the 1/9 | 1 edge half a cell right of x_e, h_s = 2h, delta/h
  % = 1/2 and 0.1; the factor (alpha_0/alpha_e)^ceil(k/2), which 2.5's
  % limitations require the caption to state; the jump's seeds the
  % monomials on the node's side, to 1e-13 in tests; first order in
  % delta/h, 2.5's table).
  \includegraphics{heat1d_stiff_seeds}
  \caption{The seeds $\phi_1, \ldots, \phi_4$ of one five-point stencil of
  the two-constant medium at $200$ nodes, whose evaluation point $x_e$ lies
  half a cell to the left of the $9 : 1$ edge (dotted), in the stencil
  coordinate $\xi = (x - x_e) / h_s$ with $h_s = 2h$; $\alpha$ across the
  stencil on top. Grey, dashed: the monomials $\xi^k$. Purple: the
  translated basis of the jump ($\delta = 0$). Blue: the seeds at
  $\delta = h/2$ (solid, with their values at the nodes) and
  $\delta = h/10$ (dashed). Each seed is multiplied by
  $(\alpha_0 / \alpha_e)^{\lceil k/2 \rceil}$, with $\alpha_0$ the jump's
  value at $x_e$ and $\alpha_e$ the edge's, so that the lines compare across
  $\delta$; the factor is the normalisation of \eqref{eq:chain}, which the
  weights never see. On the node's side of the jump its seeds are the
  monomials, and across it the slope of $\phi_1$ falls ninefold; the seeds
  of a smooth edge bend through it instead of kinking, and approach the
  translated basis at first order in $\delta / h$.}
  \label{fig:seeds1d}
\end{figure}

\para{First-order form}
% TRACE: stiff 1.3 (the state (phi_k, psi_k), psi_k = alpha phi_k'; the
% initial data; the march outward from x_e; a jump a segment boundary
% across which (phi, psi) is continuous, the weak reading; the
% double-cross), 2.3 (on each segment alpha is read from the piece of the
% medium's element that owns it, so delta = 0 is marched like any other
% width rather than dispatched to E1.2's algebra).
The chain is integrated as a first-order system in each seed and its flux
$\psi_k = \alpha \phi_k'$,
\begin{equation}
  \label{eq:firstorder}
  \phi_k' = \psi_k / \alpha, \qquad
  \psi_k' = k (k - 1)\, \alpha_e\, \phi_{k-2},
\end{equation}
with $\phi_0(x_e) = 1$, $\psi_1(x_e) = \alpha_e$ and every other initial
value zero, so that $\alpha$ is never differentiated. The ten states march
together from $x_e$ outward to each side of the stencil. Across a jump
$\phi$ and $\psi$ are continuous, which is $L$ read in the weak sense: the
march stops at the jump and reads $\alpha$ on each side from the piece
that owns it, so a jump is marched like any other width, and a stencil that
straddles both edges of a thin layer marches through both.

\para{Stencil coordinate}
% TRACE: stiff 1.3 (the chain invariant in xi; w = w~ / h_s^2), 2.3 (h_s =
% max |x_i - x_e|, xi_i = 0, +-1/2, +-1 on an interior row) and P5: 23.5
% the constant-alpha Vandermonde value, and 24 again for delta >= 10 h
% (tests, stiff 2.3); 90 at delta = 0 and 138 at delta = h/2 on the stencil
% of fig:seeds1d, paper/data/heat1d_stiff.json weights_vs_jump/matlab
% (paper_numbers.py, where "§3.1; stiff §2.3 P5").
In the coordinate $\xi = (x - x_e) / h_s$, with $h_s = \max_i |x_i - x_e|$,
the functions $\phi_k(x) / h_s^k$ satisfy the same chain with
$\alpha(x_e + h_s \xi)$ in place of $\alpha$. They are close to $\xi^k$, and
the matrix $A_{ki} = \phi_k(x_i) / h_s^k$ of \eqref{eq:moments} is a
perturbed Vandermonde matrix; the weights are $w = \tilde w / h_s^2$, with
$A \tilde w$ equal to the right-hand side of \eqref{eq:moments}. For the
centred stencil, $\xi_i = 0, \pm \tfrac12, \pm 1$, the condition number of
$A$ is $23.5$ with a constant $\alpha$. On the stencil of
Figure~\ref{fig:seeds1d}, half a cell from a $9 : 1$ edge, it is $90$ at
$\delta = 0$ and $138$ at $\delta = h/2$, and $24$ again once
$\delta \ge 10 h$, where the seeds are close to the monomials.

\para{The march and the seeded rows}
% TRACE: stiff 2.3 (DOP853 at SEED_RTOL 1e-13 / SEED_ATOL 1e-15; EDGE_STOP
% at the centre and +-10 delta; the rows whose window span meets (x_c - 20
% delta, x_c + 20 delta), TANH_REACH; 8 and 84 rows at delta = 0.0025 and
% 0.04 on 100 nodes; the batched march, some sixteen solve_ivp calls per
% operator, 0.11 s for the 1284 rows of delta = 0.04 at 1600 nodes: counts
% and timings from the notes, not a results file); 2.1 (the pieces to the
% bit beyond 20 delta); no threshold in delta, stiff 2.3's closing
% paragraph (the h <~ 2 delta regime costs the seeds nothing; at delta = 0
% they are E1.2's operator). E5.5 withdrew stiff 1.4's "an edge below
% about 1e-5 h is better served by the jump weights" (their solution is
% O(delta) from the smooth one, stiff 2.2), so it is not carried here. The
% march's floor: h^2 max |L_h u| over the seeded rows at 200 nodes,
% paper/data/heat1d_stiff.json row_residuals/matlab (paper_numbers.py,
% where "§3.1; stiff §2.3 P6"); the solution's 1.0e-11 at 101 nodes and
% delta = 5e-4 and 4e-11 at 100 and 1e-4, stiff 2.3 P6 (scratch, notes
% only); collocation not built, stiff 2.5's limitations.
The march is an explicit adaptive Runge--Kutta method of order eight
(SciPy's DOP853) at relative tolerance $10^{-13}$ and absolute tolerance
$10^{-15}$, with segment boundaries at each edge centre and at $10 \delta$
on either side of it, so that the step never has to find the edge inside a
long step over flat material. A row is seeded when the span of its stencil
meets $(x_c - 20 \delta, x_c + 20 \delta)$ for an edge centre $x_c$,
beyond which the edge is its pieces to the last bit; every other row is the
standard one. At $\delta = 0$ the seeded rows are exactly those that
straddle the jump; at $\delta = 0.0025$ and $0.04$ on $100$ nodes they
number $8$ and $84$; once $20 \delta$ covers the domain, every row is
seeded. No threshold in $\delta$ switches the seeds on or off: a resolved
edge costs them nothing (Section~\ref{sec:seedrates1d}), and at
$\delta = 0$ they are the jump's weights (Section~\ref{sec:jumplimit}).
Rows whose nodes sit at the same offsets from $x_e$ march as one system,
so an operator costs some sixteen calls of the integrator whatever the
number of seeded rows: $0.11$\,s for the $1284$ seeded rows of
$\delta = 0.04$ on $1600$ nodes. The march has one floor of its own. With
$L_h$ the discrete operator, the seeded rows' residual on the exact
equilibrium, $h^2 \max |L_h u|$, is at most $\sci{2.5}{-16}$ at $200$
nodes while the march crosses the edge in a few steps ($\delta \ge h/10$),
and $\sci{2.4}{-13}$ and $\sci{5.8}{-13}$ at $\delta = h/100$ and $h/1000$,
where it takes many. The equilibrium solution shows it once
$\delta \lesssim h/40$, at $10^{-11}$ to $\sci{4}{-11}$. A collocation
integration of the same chain would remove the floor; we have not built
one.

\subsection{Nested kernels and the jump as the limit}
\label{sec:jumplimit}
% TRACE: stiff 1.4 (the nested spaces; the spaces, not the basis,
% determine the weights; the weak reading; the translated x with slopes in
% the ratio alpha^+/alpha^-; for constant pieces the chain is the
% translated basis, one construction with two implementations); 2.3 P4 and
% 2.5 (the march against E1.2's algebra at delta = 0: 7e-16 in the
% weights, 3e-15 in the functions, paper/data/heat1d_stiff.json
% weights_vs_jump/matlab and seed_functions, paper_numbers.py where "§3.2;
% stiff §2.5 (2)"); LITERATURE.md 1a K3 (Sujecki, Gartland).
Since $L$ sends each seed two steps down the chain,
$\operatorname{span}\{\phi_0, \ldots, \phi_k\}$ is the $(k + 1)$-th space
of
\begin{equation}
  \label{eq:nested}
  \{1\} \subset \ker L \subset \ker \px L \subset \ker L^2 \subset
  \ker \px L^2 \subset \cdots,
\end{equation}
where $\ker \px L^m$ holds the $g$ with $L^m g$ constant. The spaces depend
on neither the anchor nor the normalisation, so the weights do not either;
only the basis does.

For a jump, read $L$ in the weak sense, $\phi$ and $\alpha \phi'$
continuous across it. Then $\ker L$ is spanned by $1$ and the
piecewise-linear function whose slope changes by the factor
$\alpha^- / \alpha^+$ across the jump, the translated $x$ of
Section~\ref{sec:jump}. On constant pieces $L \phi_2 = 2 \alpha_e$ makes
$\phi_2'' = 2 \alpha_e / \alpha^\pm$ on the two sides, with $\phi_2$ and its
flux continuous, which is the continuity of $u_t$ and of its flux; and so
on up the chain, each condition $\pt^m u = \text{const}$ or
$\alpha\, \px \pt^m u = \text{const}$ being one of \eqref{eq:continuity}.
For constant pieces the chain across a jump is therefore the translated
basis, the anchor conditions picking one basis of its span, and the
interface-aware stencils are the seed stencils at $\delta = 0$: one
construction with two implementations, algebra for a jump and an ODE march
for a smooth edge. On the stencil of Figure~\ref{fig:seeds1d} the march
reproduces the translated basis to $\sci{3}{-15}$ function by function, and
its weights to $\sci{7}{-16}$. The same device, Taylor data carried across
a step to make the weights, appears in photonics \cite{Sujecki2010}, and
three-point schemes of any prescribed order across a jump in the
coefficient are Gartland's \cite{Gartland1989}.

% TRACE: stiff 2.3 P4 (0.843, 0.465, 0.114, 0.0107, 0.00107 of the largest
% weight at delta/h = 1, 1/2, 0.1, 0.01, 0.001), paper/data/heat1d_stiff.json
% weights_vs_jump/matlab (paper_numbers.py, where "§3.2; stiff §2.3 P4");
% E5.5's number check: 0.465 is 46 %, not the 47 % of stiff 1.4's scratch;
% the functions first order too, stiff 2.5's table (5.12e-2, 5.12e-3,
% 5.12e-4 for phi_1 at delta/h = 0.1, 0.01, 0.001, tests); P8, the
% double-cross, stiff 2.3 (1 | 0.1 | 1 on [0, 2h] at h = 0.05: 7.6e-15 at
% delta = 0; 22.2 %, 1.86 %, 0.182 % at delta/h = 0.1, 0.01, 0.001; tests).
As $\delta \to 0$ the seed weights approach the jump's at first order in
$\delta / h$. On the stencil of Figure~\ref{fig:seeds1d},
$\max_i |w_i - w_i^{\mathrm{jump}}| / \max_i |w_i^{\mathrm{jump}}|$ is
$84$, $46$, $11$, $1.1$ and $0.11$ percent at $\delta / h = 1$,
$\tfrac12$, $0.1$, $0.01$ and $0.001$, and the seeds approach the
translated basis at the same rate. A stencil that straddles both edges of a
layer two cells thick, $\alpha = 1 \,|\, 0.1 \,|\, 1$, marches through
both: at $\delta = 0$ its weights are those of the jump construction
translated twice, to $\sci{7.6}{-15}$, and at $\delta / h = 0.1$, $0.01$
and $0.001$ they differ from them by $22$, $1.9$ and $0.18$ percent.

% TRACE: stiff 1.4 (phi_2 = (x - x_e)^2 - (2/3)(alpha_e'/alpha_e)(x -
% x_e)^3 + ...; the two spaces differ at relative O(h alpha'/alpha); both
% exact on their own space), 2.3 P4 on eq. 75 (the row whose evaluation
% point is the layer's edge x = 0, 201 nodes, alpha_e = 0.1 the layer's,
% alpha'/alpha = 25: 1.25 at delta = 0 and 1.26-2.24 for delta/h from 0.001
% to 1, paper/data/heat1d_stiff.json weights_vs_jump/eq75, paper_numbers.py
% where "§3.2; stiff §2.3 P4"; first order in h at x = 0.5, 1.75, 1.25,
% 0.50, 0.22 for h = 0.02 to 0.0025, tests), 2.5 (the functions 0.19-0.40
% apart, so eq. 75's seeds are read against the reference solution); the
% seed operator better by 6-44 times at equilibrium and 6-38 on the ramp
% problem at delta = 0 over 101-1601 nodes, knee/equilibrium and knee/ramp
% eq75 at delta = 0, where the construction's column is E1.2's operator
% (paper_numbers.py, where "§3.2; stiff §2.3 P6" and "P7").
When the pieces vary the limit is not exact. The interface-aware
construction expands each side's $\alpha$ to degree four about the
interface and truncates the operator there (Section~\ref{sec:jump}), while
the seed $\phi_2$ solves $(\alpha \phi_2')' = 2 \alpha_e$ exactly,
$\phi_2 = (x - x_e)^2 - \tfrac23\, (\alpha'(x_e) / \alpha_e)\, (x - x_e)^3
+ \cdots$, so that within a stencil the two five-dimensional spaces differ
at relative order $h \alpha' / \alpha$. On the sinusoidal medium, whose
layer has $\alpha' / \alpha = 25$ at its edge, the seed and jump weights
of the row whose evaluation point is that edge differ at $h = 0.01$ by
$1.25$ times the largest weight at $\delta = 0$ and by $1.26$ to $2.24$
times for $\delta / h$ from $0.001$ to $1$, and the difference falls with
$h$.
Neither space is wrong, each being exact on its own five functions; on this
medium the seed operator's error is the smaller, by $6$ to $44$ times at
equilibrium and $6$ to $38$ times on the ramp problem at $\delta = 0$ over
$101$ to $1601$ nodes. There the seeds are measured against the reference
solution, never against the jump construction.

\subsection{Well-posedness}
\label{sec:ect}
% TRACE: stiff 1.5 (the proposition and its Polya form; the owner's
% alpha_e on a jump) with E5.5's note there: a complete Chebyshev system
% for piecewise-continuous alpha, an extended one where alpha is smooth
% (the classical theorem asks the weights for derivatives), proved from
% eq:nestedint, the seeds being multiples of the canonical nested
% integrals; LITERATURE.md 1a K7 (the companion's argument, on Polya,
% Karlin-Studden, Coppel).
The stencil solve is nonsingular for every edge width, the jump included,
and for every set of distinct nodes, not only the ones we used.

\begin{proposition}
\label{prop:ect}
Let $\alpha$ be positive and bounded away from zero on an interval $I$
containing $x_e$, and continuous on $I$ except at finitely many jumps,
across which $L$ is read in the weak sense. Then for every $k$ the seeds
$\phi_0, \ldots, \phi_k$ form a complete Chebyshev system on $I$:
interpolation in their span at any $k + 1$ distinct points of $I$ is
uniquely solvable. Where $\alpha$ is smooth, as through a $\tanh$ edge
between smooth pieces, they form an extended complete Chebyshev system. In
particular \eqref{eq:moments} has a unique solution for every
$\delta \ge 0$ and every five distinct nodes in $I$.
\end{proposition}

\begin{proof}
By \eqref{eq:nestedint} the seeds are nonzero multiples of $I_0, \ldots,
I_k$, nested integrals whose weights $v_1 = 1/\alpha$, $v_2 = 1$,
$v_3 = 1/\alpha, \ldots$ (outermost first) are positive. Such a family is
the canonical basis of the kernel of the operator in P\'olya form
$\px \frac{1}{v_k} \px \cdots \frac{1}{v_2} \px \frac{1}{v_1} \px$, which is
$\px L^m$ for $k = 2m$ and $L^{m+1}$ for $k = 2m + 1$, and whose kernel
is the $(k + 1)$-th space of \eqref{eq:nested}. Suppose $g = \sum_{j \le k} c_j I_j$
vanishes at $k + 1$ distinct points of $I$. Then $g' = v_1 g_1$, where
$g_1 = \sum_{1 \le j \le k} c_j I'_{j-1}$ is a combination of the $k$
nested integrals $I'_0, \ldots, I'_{k-1}$ with the weights $v_2, v_3,
\ldots$, and $g_1$ is continuous. Between two consecutive zeros $a < b$ of
$g$, $\int_a^b v_1 g_1 = 0$ with $v_1 > 0$, so $g_1$ vanishes in $(a, b)$:
it has $k$ distinct zeros. By induction on $k$, $g_1 \equiv 0$, so
$c_1 = \cdots = c_k = 0$, and then $g = c_0$ vanishes, so $c_0 = 0$. Only
the positivity and integrability of the weights entered, so the argument
holds across finitely many jumps. Where $\alpha$ is smooth the operator is
disconjugate on $I$ and the family is an extended complete Chebyshev system
\cite{Polya1922,KarlinStudden1966,Coppel1971}.
\end{proof}

% TRACE: stiff 1.5 (x_e on a jump, as for E1.2's rows at an interface
% node: alpha_e the owner's value, region_index putting a point on a jump
% to its right, PiecewiseAlpha.at_interface naming the owner; phi_1 slope
% 1 on the owner's side and alpha_e/alpha_other on the other); LITERATURE.md
% 1a K7 (the argument is the companion's).
When $x_e$ lies on a jump, as it does for the rows whose evaluation node
falls on an edge of the sinusoidal medium, $\alpha_e$ is the value of the
piece that owns $x_e$; $\phi_1$ then has slope $1$ on that side and
$\alpha_e / \alpha$ on the other, and nothing else changes, since
$\alpha_e$ enters the seeds only as a factor. The argument is the one the
companion paper makes for the two fields of the wave equation
\cite{Martin2026seeds}, shortened by \eqref{eq:nestedint}.

\subsection{Truncation error}
\label{sec:truncation}
% TRACE: stiff 1.6 (the Newton form with Muhlbach's remainder; M = d_x L^2,
% five derivatives; tau(g) = c_5 h^3 (M g) + ...; one order below the
% companion's first-derivative bound; E1.2's h^3 moments 0.32, 0.084 and
% 172 for the 1/9 | 1 jump at 1/2, 3/2 and 0 cells, zero to rounding for
% equal pieces, tests/heat1d/test_interface.py::test_jump_rows_have_an_h3_
% moment_and_equal_pieces_do_not); the global order four, stiff 1.6 (the
% bounded discrete Green's function; E1's rates 3.85-4.19 on the jump,
% port notes 1.3-1.4); LITERATURE.md 1a K7 (Beale-Layton's O(h)
% truncation at the interface, O(h^2) solution: cited here, not claimed).
Interpolation in an extended complete Chebyshev system has a Newton form
whose remainder carries the operator that annihilates the space where the
polynomial remainder carries the $(k + 1)$-th derivative
\cite{Muhlbach1973}. For the five seeds that operator is $M = \px L^2$, and
the local truncation error of a seed stencil on a smooth function $g$ is
\begin{equation}
  \label{eq:truncation}
  \tau(g) = \sum_{i=1}^{5} w_i\, g(x_i) - (L g)(x_e) = O(h^3\, M g),
\end{equation}
with a constant that the remainder expresses through $1/\alpha$ inside the
stencil. That is third order, not fourth: a second derivative on five
points is generically third order, and the centred standard stencil gains
its fourth only because its $h^3$ term vanishes by symmetry, which an edge
breaks. The interface-aware rows of a jump show the same. On the function
of degree five in their translated basis, which lies outside the space they
reproduce, the $h^3$ term is $0.32$, $0.084$ and $172$ (in the stencil
coordinate) for a $9 : 1$ jump half a cell and one and a half cells from
the evaluation node and at the node itself, and zero to rounding for equal
pieces.

The global order is still four. A local error of order $h^3$ on a fixed
number of rows costs one order less in the solution than the same error on
every row, since the inverse of the discrete operator has entries of
order $h$, so the solution error is $O(h^4)$. This is what the
interface-aware stencils reach at a jump (rates $3.85$ to $4.19$,
elliptic and parabolic, in our port of \cite{Martin2016}), and it is the
phenomenon Beale and Layton prove for interface corrections, where an
$O(h)$ truncation error at the interface leaves the solution second order
\cite{BealeLayton2006}. The bound \eqref{eq:truncation} has two readings,
and Section~\ref{sec:results1d} is built on them.

\para{Equilibrium}
% TRACE: stiff 1.6's first reading (u = A + B phi_1 in every seeded
% window; the seeded rows exact at any delta and h).
If $L u = 0$ then $M u = 0$, and more than that: $u$ has a constant flux,
so on every seeded stencil it lies in $\operatorname{span}\{\phi_0,
\phi_1\}$, and \eqref{eq:moments} reproduces $L u$ there exactly, at any
$\delta$ and any $h$ (Section~\ref{sec:elliptic1d}).

\para{Time dependence}
% TRACE: stiff 1.6's second reading (M u = d_x u_tt, bounded independently
% of delta, changing by the contrast ratio across the edge; FD4's u^(5) ~
% delta^-4; the naive rows miss an O(1) fraction of the slope jump, O(1/h)
% in L u; the constant's delta-independence not a consequence of the
% proposition, P7); 2.3 (the separable solution e^{ct} v on the
% two-constant medium, 100 -> 400 nodes: the seeded rows' residual rates
% 3.06, 3.78 at delta = 0.01 and 3.15, 3.51 at 0.0025, residual n^3 0.27,
% 0.26, 0.15 and 0.43, 0.38, 0.27, within a factor two between the widths;
% the plain rows' 3.94, 3.97 and 3.94, 4.01; tests, notes only); the
% smoothing paragraph, stiff 1.6's last (the flux changes by int u_t dx =
% O(delta) across the edge; why the construction is a serious baseline).
If $u_t = L u$ then $L^2 u = u_{tt}$ and $M u = \px u_{tt}$, the
$x$-derivative of another solution of the heat equation. Since $u_{tt}$
has a continuous flux, $\px u_{tt}$ is bounded independently of $\delta$
and changes across the edge by the ratio of the pieces, whatever the width.
The standard stencil's error is driven by $u^{(5)}$ instead, which grows
like $\delta^{-4}$ inside the edge, and the rows of $D_x A D_x$ that see an
unresolved edge miss a fixed fraction of the jump in slope, an error of
order $1/h$ in $L u$. The bound does not make the constant in
\eqref{eq:truncation} independent of $\delta / h$, since the remainder
involves $1/\alpha$ across the stencil. Measured on a separable solution
$e^{ct} v(x)$ of the two-constant medium from $100$ to $400$ nodes, the
seeded rows' residual is third order (rates $3.06$ to $3.78$), its
constants at $\delta = 0.01$ and $0.0025$ within a factor of two of each
other, and the standard rows' residual is fourth order ($3.94$ to
$4.01$).

Diffusion smooths, so for $t > 0$ every $\pt^k u$ is itself a solution and
the second reading holds for the whole chain. And since
$\px (\alpha u_x) = u_t$, the flux changes across an edge of width $\delta$
only by $\int u_t\, dx = O(\delta)$: in diffusion a sub-grid edge is close
to a jump at all times. That is why the $\delta = 0$ construction is a
serious baseline here (Section~\ref{sec:construction1d}), and why the
seeds' margin over it is set by $\delta$.

\subsection{The elliptic problem in one dimension}
\label{sec:elliptic1d}
% TRACE: stiff 1.8 (the exact constructions; the face conductance a_{i+1/2}
% = h / int dxi / alpha, exact for any profile; the plan's T1 not item 3;
% the correction for E3.5 with E5.5's note there: no cell mean sampled at
% the nodes makes Dx A Dx exact, while a nodal alpha read off the discrete
% solution would; tests/heat1d/test_operators.py::test_no_nodal_alpha_
% makes_dx_a_dx_exact_on_a_kinked_equilibrium, a residual above 0.5 at 100
% nodes for one and two cells), 1.6's first reading, 2.5 (4) (the
% sinusoidal medium: the plain rows not exact, the conservative scheme
% exact); LITERATURE.md 1a K1 (Tikhonov-Samarskii, Samarskii; Patankar
% section 4.2-3, pp. 44-47, eqs. 4.9-4.10, derived from the steady,
% no-source, one-dimensional case; Pan-Xu-Li's harmonic average method,
% their eq. 3, second order their theorem; Goode-Appel's statement of
% exactness), K10 (i) (T1 the like-for-like changed-medium comparator,
% citing nothing); plan 3.2 and R2.
The one-dimensional equilibrium has a constant flux $\alpha u' = B$, so
$u = u(-1) + B \int_{-1}^{x} d\xi / \alpha$, and two constructions are
exact on it at every $h$ and every $\delta$ when the pieces are constant.

The first is the seed operator. Its seeded rows are exact by the first
reading of Section~\ref{sec:truncation}, and its standard rows see a
linear function on a constant piece. At $\delta = 0$ it is the
interface-aware operator of the jump.

The second is the conservative three-point scheme with the face
conductances
\begin{equation}
  \label{eq:faceconductance}
  a_{i+1/2} = h \Big/ \int_{x_i}^{x_{i+1}} \frac{d\xi}{\alpha(\xi)},
\end{equation}
whose discrete flux $a_{i+1/2} (u_{i+1} - u_i) / h$ equals $B$ at the
exact nodal values whatever $\alpha$ is, so that the exact solution
satisfies every discrete equation. It is the homogeneous scheme of
Tikhonov and Samarskii \cite{TikhonovSamarskii1962,Samarskii2001}, whose
face coefficient is the conductance of the constant-flux seed over one
cell. For a piecewise-constant $\alpha$ it is the harmonic-mean interface
conductivity of finite-volume heat conduction, which Patankar derives from
the same steady one-dimensional flux \cite[\S4.2-3]{Patankar1980} and which
is the exact interblock transmissivity of groundwater models for steady
one-dimensional flow through blocks of uniform transmissivity
\cite{GoodeAppel1992}. Pan, Xu and Li call the scheme
the harmonic average method and prove it second order
\cite{PanXuLi2026}. The exactness belongs to the conservative form: the
same harmonic mean sampled at the nodes and used as $A$ in $D_x A D_x$,
the like-for-like changed medium of Section~\ref{sec:treatments1d}, leaves
a residual above $0.5$ on the exact equilibrium of the $9 : 1$ jump at
$100$ nodes, whether the mean is taken over one cell or two.

So on constant pieces the one-dimensional equilibrium cannot rank the seeds
against the conservative scheme. It is the sanity check and the limit
case, and the one-dimensional ranking is the parabolic problem's, where
the conservative scheme is second order and the seeds are fourth
(Section~\ref{sec:treatments1d}). Both exactness results are classical,
and so is the conservative scheme's second order; we observe them and claim
none of them. Where the pieces vary, as in the sinusoidal medium, the
equilibrium does rank the two, for a reason that has nothing to do with
the edge: the conservative scheme stays exact, while the seed operator's
standard rows inside the layer keep the error of a smooth coefficient, and
it is exact only where the seeded rows cover the layer
(Section~\ref{sec:seedrates1d}).

% ===========================================================================
\section{One-dimensional results}
\label{sec:results1d}
% TRACE: stiff 2.1 (media and references), 2.2 (the knee, the delta = 0
% construction), 2.3 (the seeds), 2.4 (the treatments), 2.5 (the four
% statements, the snapshot); numbers from paper/data/heat1d_stiff.json
% (#44), asserted by scripts/paper_numbers.py. Drafted by #47.
\subsection{Media and references}
\label{sec:media1d}
\stub{\#47: the smooth-edged media and the references at every $\delta$}

\subsection{The parabolic knee}
\label{sec:knee1d}
\stub{\#47: the naive knee at $h \approx \delta$; the snapshot of the four
regimes on one grid}

\subsection{The jump construction on a smooth edge}
\label{sec:construction1d}
\stub{\#47: the $\delta = 0$ construction's error on a smooth edge}

\subsection{The seeds}
\label{sec:seedrates1d}
\stub{\#47: the seeds' rates at every $\delta$}

\subsection{Coefficient treatments}
\label{sec:treatments1d}
\stub{\#47: the treatments' table and the ranking; the elliptic sanity
figure as a remark}

% ===========================================================================
\section{Seeds in two dimensions}
\label{sec:seeds2d}
% TRACE: stiff 3.2-3.4 (the straight feature), 3.10 (the tangential chain),
% 3.11 (the ring's gap, series, flux seeds and level-0 warp), 4.10 and 5.2
% (the diagonal rule, and why it is the ring's); code heat2d/seeds.py.
% Drafted by #48. Figure: a seed section (#44).
\subsection{A straight feature}
\label{sec:straight}
\stub{\#48: the local frame and the anchor; seeds polynomial in the
tangential coordinate with coefficient functions marched along the normal}

\subsection{Stencil assembly}
\label{sec:assembly2d}
\stub{\#48: the seeds in the RBF-FD system, the warp they bring with them,
which rows are seeded; the seed-section figure}

\subsection{Curved and tangentially varying features}
\label{sec:curved}
\stub{\#48: why a frozen normal profile is inconsistent where $\alpha$
varies along the edge; the tangential chain in the curve's own coordinates}

\subsection{A thin resistive layer}
\label{sec:ring}
% For #48, moved here by #45 from LITERATURE.md 6a (its third cut; section 1
% keeps one clause, "A resistive layer thinner than the grid is an imperfect
% interface in the limit [...], which Section 5.4 takes up"). Place it where
% the flux seeds are introduced, with 6b's required flux-seed sentence:
%   A resistive layer thinner than the grid is an imperfect interface in the
%   limit \cite{SanchezPalencia1974,PerrusselPoignard2013,Hashin2001}, and
%   strong-form schemes across such a contact need its jump conditions one
%   derivative higher in two dimensions \cite{LombardPiraux2003,LombardPiraux2006};
%   the flux seeds of Section~\ref{sec:ring} are that observation for the
%   diffusion operator.
\stub{\#48: the double edge of a ring: its exact width, series composition,
the flux seeds, the warp on rows anchored inside the layer}

% ===========================================================================
\section{Two-dimensional results}
\label{sec:results2d}
% TRACE: stiff 4.1-4.10 and 5.1 (the snapshot, the results files), in the
% order of stiff 5.7; the statement numbers are stiff 5.3's; numbers from
% paper/data/ (#44). Drafted by #49. Every table from a figures/tab_*.tex
% fragment.
\subsection{Test problems and references}
\label{sec:problems2d}
\stub{\#49: the flat band, the sine pair, the ring; the separable, product-grid
and reference-free references (stiff 4.1, 4.6, 4.8; statement 1)}

\subsection{Naive baseline}
\label{sec:naive2d}
\stub{\#49: the knee on scattered nodes and the jump construction on its
floor (stiff 4.2; statements 2--3)}

\subsection{Solvability and spectra}
\label{sec:spectra2d}
\stub{\#49: the seed rows in the global matrix (stiff 4.4; statement 6)}

\subsection{The flat sweep}
\label{sec:flat2d}
\stub{\#49: the seeds against the baselines at every $(n, \delta)$, the
warp, the elliptic ranking (stiff 4.5; statements 4, 5, 7)}

\subsection{The curved feature}
\label{sec:curved2d}
\stub{\#49: the frozen profile and the tangential chain on the sine pair
(stiff 4.6--4.7; statement 8)}

\subsection{The ring and the extremizing sweep}
\label{sec:ring2d}
\stub{\#49: the twins of EABE Fig.~19--20 and the smooth ring (stiff 4.8,
4.10; statement 9)}

\subsection{Coefficient treatments}
\label{sec:treatments2d}
\stub{\#49: the treatments on scattered nodes (stiff 4.9; statement 10)}

\subsection{A snapshot}
\label{sec:snapshot2d}
\stub{\#49: the four regimes on one node set (stiff 5.1; statement 11)}

% ===========================================================================
\section{Discussion}
\label{sec:discussion}
% TRACE: stiff 2.5 (the 1-D limitations), 5.5 (limitations), 5.6 (the
% roundtable's deferred runs); plan D10 (the corner); LITERATURE.md 6b
% (#43) for the conclusions, verbatim. Drafted by #50.
\subsection{Limitations}
\label{sec:limitations}
\stub{\#50: corners, extreme contrasts, the cost in three dimensions, what
the experiments leave open}

\subsection{Future work}
\label{sec:future}
\stub{\#50: in the order to do it}

\subsection{Conclusions}
\label{sec:conclusions}
\stub{\#50: LITERATURE.md section 6b, verbatim}

\section*{Acknowledgments}
% Drafted by #52.
\stub{\#52: acknowledgments; where the code and data live}

\section*{Tool and computational resource disclosure}
% Drafted by #52, in the form used for the companion manuscript.
\stub{\#52}

% The scaffold cites nothing yet; \nocite{*} keeps the seeded bibliography
% building until the drafts cite it. #51 removes it, and make_arxiv.py
% refuses to package while it is here.
\nocite{*}
\bibliographystyle{amsalpha}
\bibliography{references}

\end{document}
